\section{Integral coordinates and universal classes}\label{sec:ring}

We now give the integral coordinates in which the compact wall and
attachment calculations can be performed. The Hilbert-square basis
and Markman transport are classical inputs. The purpose of this
section is to fix their signs and the universal normalization for
this moduli problem, including arbitrary topological index classes.

\subsection{The diagonal construction}
Let $\alpha_1,\ldots,\alpha_{22}$ be an integral basis of $\Lambda$,
write $q_{ij}=\alpha_i\cdot\alpha_j$, and let $(q^{ij})$ be its
inverse matrix. Let $x_X$ be the oriented point class of $X$.
Consider
\[
 b:B=\Bl_\Delta(X\times X)\longrightarrow X\times X,
 \qquad p:B\longrightarrow Y=X^{[2]}.
\]
The map $p$ is the double cover ramified along the exceptional
divisor $E_B$. Write $j:E_B\to B$, $\pi:E_B\to X$ and
$e=[E_B]=p^*\delta$. A tensor class below is pulled back by $b$,
and $e\alpha$ means $j_*\pi^*\alpha$. The blow-up relations are
\begin{equation}\label{eq:blowup}
 \begin{gathered}
 (e\alpha)(u\otimes v)=e(\alpha uv),\qquad
 (e\alpha)(e\beta)=-[\Delta](\alpha\beta\otimes1),\\
 [\Delta]=x_X\otimes1+1\otimes x_X+
             \sum_{i,j}q^{ij}\alpha_i\otimes\alpha_j.
 \end{gathered}
\end{equation}
For the sign, $E_B|_{E_B}=-\xi$ and the excess formula gives
$b^*[\Delta]=j_*\xi$, since $c_1(TX)=0$. Hence $e^2=-b^*[\Delta]$.
In particular $e^4=24x_X\otimes x_X$. Integration on $Y$ is half
integration of a pullback on $B$; $p^*[\pt]=2x_X\otimes x_X$.
Thus \eqref{eq:blowup} is a multiplication rule using only the
marked K3 pairing.

The pullback of $\mu(\alpha)$ is $\alpha\otimes1+1\otimes\alpha$.
The classes $\mu_i=\mu(\alpha_i)$ and $\delta$ form an integral
basis in degree two. We extend $q$ to their Beauville--Bogomolov
pairing by
\begin{equation}\label{eq:BB}
 q(\mu(a),\mu(b))=q(a,b),\quad q(\mu(a),\delta)=0,\quad
 q(\delta,\delta)=-2.
\end{equation}
An integral degree-four basis consists of the classes whose
pullbacks are
\begin{equation}\label{eq:basis4}
 \begin{aligned}
 p^*A&=x_X\otimes1+1\otimes x_X,\qquad &p^*T_i&=2e\alpha_i,\\
 p^*C_{ij}&=\alpha_i\otimes\alpha_j+\alpha_j\otimes\alpha_i
       \quad(i<j),&p^*Q_i&=\alpha_i\otimes\alpha_i-e\alpha_i.
 \end{aligned}
\end{equation}
Integrality is Qin--Wang's theorem, with its K3 extension
\cite[Theorem~5.4 and Remark~5.6]{qw}, in the geometric
Nakajima conventions of \cite[Remark~5.7]{bns}. In particular
$Q_i$ is the half difference of the two Nakajima terms, not
half a symmetric square by itself. These statements apply to
all of $\Lambda$, including its nonalgebraic classes.
We use the basis statement in that reference; the product
signs here follow directly from \eqref{eq:blowup}.

For completeness we print the entire middle pairing. Degree-eight
products in the following display denote their integral against $Y$:
\begin{equation}\label{eq:gram}
 \begin{aligned}
 A^2&=1,& AT_i=AC_{ij}=AQ_i&=0,\\
 T_iT_j&=-2q_{ij},& T_iQ_j&=q_{ij},&T_iC_{jk}&=0,\\
 Q_iQ_j&=(q_{ij}^2-q_{ij})/2,&&&Q_iC_{jk}&=q_{ij}q_{ik},\\
 C_{ij}C_{kl}&=q_{ik}q_{jl}+q_{il}q_{jk}.
 \end{aligned}
\end{equation}
Every entry follows by multiplication on $B$ and division by two.
The pairing is unimodular. Eliminate the $T$ block $-2q$ over
$\Q$. Its Schur complement is half the tensor-product pairing
on the symmetric tensors in $\Lambda\otimes\Lambda$; its
determinant is $2^{-22}(\det q)^{23}$. To verify this identity,
diagonalize $q$ over a field extension: the diagonal symmetric
tensors have entries $d_i^2/2$, and the off-diagonal ones have
entries $d_i d_j$. The product is the stated expression, and
the identity descends over $\Q$. The determinant of the whole
pairing is therefore
$\det(-2q)2^{-22}(\det q)^{23}=(\det q)^{24}=1$.
This also verifies index one, once integrality and the known
free rank $276$ are supplied.

Products of degree-two classes are
\begin{equation}\label{eq:degree2-products}
 \begin{aligned}
 \mu_i\mu_j&=q_{ij}A+C_{ij} &&(i<j),\\
 \mu_i^2&=q_{ii}A+2Q_i+T_i,&\mu_i\delta&=T_i,\\
 \delta^2&=-A-\sum_iq^{ii}Q_i-\tfrac12\sum_iq^{ii}T_i
                     -\sum_{i<j}q^{ij}C_{ij}.
 \end{aligned}
\end{equation}
The diagonal of $q^{-1}$ is even because the dual K3 lattice is
even. Thus the last expression is integral. If $\alpha^i$ is the
$q$-dual basis, set
\begin{equation}\label{eq:basis6}
 R_i=\mu(\alpha^i)A,\qquad R_\delta=-\delta A.
\end{equation}
These classes are dual to $\mu_i,\delta$ and form an integral
degree-six basis. For $a\in R^2$ and $z\in R^4$, its product
$az$ has coefficient $\int_Y\mu_i az$ on $R_i$ and
$\int_Y\delta az$ on $R_\delta$. These are given by
\eqref{eq:degree2-products} and \eqref{eq:gram}.
Together with the unit, the oriented point, the dual degree-two/
degree-six pairing, and zero odd groups, this specifies every product.
The ranks are $1,23,276,23,1$. Torsion-freeness is the integral
Hilbert-scheme theorem; see \cite[Theorem~3.1]{totaro}.

The formulas also exhibit rational generation by degree two:
\begin{equation}\label{eq:rational-generation}
 A=\frac1{10}\delta^2+\frac1{20}\sum_{i,j}q^{ij}\mu_i\mu_j,
\end{equation}
after which \eqref{eq:degree2-products} gives $T_i,C_{ij},Q_i$,
and \eqref{eq:basis6} gives degree six. The integral ring generally
requires more than the integral symmetric square.

We will also use the classical tangent-class identity
\begin{equation}\label{eq:c2}
 c_2(Y)=24A-3\delta^2
\end{equation}
of Shen--Vial \cite[Lemma~14.4]{sv}. Their incidence-surface
class is $A$ here. Their rational Chow identity gives this integral
cohomological identity by torsion-freeness. The preceding products
give, independently,
\[
 \int_Y\delta^4=12,\quad
 \int_Y c_2\mu(a)\mu(b)=30q(a,b),\quad
 \int_Yc_2\delta^2=-60,\quad \int_Yc_2^2=828.
\]

\subsection{Transport from the rank-three moduli}
Set $a_H=(1,H,2)$, $C_0=D_1+D_2-2H$ and
\begin{equation}\label{eq:g}
 g=e^{-C_0}s_{a_H},\qquad s_{a_H}(z)=z+(z,a_H)a_H.
\end{equation}
Here $e^D(r,c,s)=(r,c+rD,s+cD+rD^2/2)$.
The input lattice gives
\[
 a_H^2=-2,\quad (v,a_H)=-2,\quad C_0^2=-14,\quad
 s_{a_H}(v)=(1,C_0,-8),\quad g(v)=v'=(1,0,-1).
\]
This is the integral Mukai isometry induced by the spherical twist
about $\OO_X(H)$ followed by tensoring with $\OO_X(-C_0)$.
It need not send each chamber sheaf to a sheaf.

The polarizations may be chosen in each sector with degree coprime
to three, as in \cite[\S4.5]{wuebben17}; the full compact
Gieseker moduli therefore satisfies Markman's suitability condition
\cite[Definition~3.4]{markman08}. His universal Chern
correspondence \cite[Theorem~1.2 and equation~(3)]{markman08}
gives
\begin{equation}\label{eq:Phi}
 \Phi_\sigma=\gamma_{g,v}:H^*(M_\sigma,\Z)\xrightarrow{\sim} R.
\end{equation}
The source and target are torsion-free: for the source use the
K3$^{[2]}$ deformation type established in Section~\ref{sec:geometry},
and for the Hilbert square use \cite[Theorem~3.1]{totaro}.
Thus the modulo-torsion qualification loses nothing here.
The correspondence preserves tangent Chern classes and
orientation; for the latter, preservation of $c_4=324[\pt]$
suffices. Compatibility with the Mukai homomorphisms is
\cite[Theorem~3.9 and diagram~(41)]{markman08}:
\begin{equation}\label{eq:theta-transport}
 \Phi_\sigma\theta_v(z)=\theta_{v'}(gz),\qquad
 \theta_{v'}(r,c,r)=\mu(c)+r\delta.
\end{equation}
The sign in the second formula follows from the universal ideal
sequence: $c_1(\OO_X^{[2]})=-\delta$, while the index of
$(r,-c,r)$ is its Euler characteristic minus its tautological
class. This determines $\Phi_\sigma$ without choosing an
unspecified deformation path.

\subsection{Topological indices and determinant normalization}
Let $\mathcal I$ be the universal ideal class on $X\times Y$.
For a topological class $F$ let $F^{[2]}$ denote the index
defined by the universal length-two subscheme. With $g$ acting
on topological classes through their Mukai vectors, put
\begin{equation}\label{eq:J}
 J_z=\chi_X((gz)^\vee)\,1-((gz)^\vee)^{[2]},\qquad
 \mathcal L=\det J_y,\qquad \widehat J_z=J_z\otimes\mathcal L^{-1}.
\end{equation}
Here $gy=(0,-f_2,1)$ and $\rank J_y=1$.

\begin{proposition}\label{prop:indices}
For every $z\in K^0_{\mathrm{top}}(X)$ and $k\ge0$,
\begin{equation}\label{eq:indices-transport}
 \Phi_\sigma(c_k(I_z^y))=c_k(\widehat J_z).
\end{equation}
These classes are calculated in the integral basis above by
\begin{equation}\label{eq:tautological-character}
 p^*\ch(F^{[2]})=\ch(F)_1e^{-e}+\ch(F)_2,
 \qquad F\in K^0_{\mathrm{top}}(X).
\end{equation}
Subscripts denote the two pullbacks from $X$ to $B$.
\end{proposition}

\begin{proof}
Markman's universal-class identity intertwines the universal Chern
characters up to a topological base line bundle. Contract its K3
factor against $gz$ using the Euler pairing. Since $g$ is an
Euler isometry, the index characters are intertwined up to that
same line. Taking the determinant of the rank-one $y$-index
multiplies it by exactly this line, so normalization cancels it.
The resulting normalized characters agree over $\Q$, and hence
so do the Chern classes. Both Chern classes are integral and
$H^*(Y,\Z)$ is torsion-free, proving
\eqref{eq:indices-transport} integrally.

We prove \eqref{eq:tautological-character} for topological $F$,
without requiring an algebraic representative of its first
Chern class. Let $\mathcal Z\subset Y\times X$ be the
universal subscheme and pull it back to $\mathcal Z_B\subset
B\times X$. The incidence description is the diagonal blow-up
\cite[\S10, diagram~(54)]{sv}. If $q_1,q_2:B\to X$
are the projections, $\mathcal Z_B$ is the scheme-theoretic
union of their graphs, intersecting in the graph over $E_B$.
This includes the nonreduced fibre: in a blow-up chart with
difference coordinates $(t,tw)$ the length-two algebra is
$A[U]/(U(U-t))$. Evaluation at $0,t$ maps it injectively to
$A\oplus A$, with image the pairs agreeing modulo $t$ and
cokernel $A/(t)$. The other blow-up chart gives the same
description. Thus there is an exact sequence of kernels
\begin{equation}\label{eq:kernel}
 0\longrightarrow\OO_{\mathcal Z_B}\longrightarrow
 \OO_{\Gamma_1}\oplus\OO_{\Gamma_2}\longrightarrow
 \OO_{\Gamma_{E_B}}\longrightarrow0.
\end{equation}
All these kernels are perfect on the smooth product $B\times X$.
Pass their equality to topological $K$-theory, tensor with the
pullback of any $F$, and push forward to $B$. Base change for
the complex-oriented projection and functoriality give
\[
 p^*F^{[2]}=q_1^*F+q_2^*F-j_!\pi^*F
           =q_1^*F\otimes\OO_B(-E_B)+q_2^*F.
\]
The last equality uses the divisor Koszul class
$j_!1=1-[\OO_B(-E_B)]$ and the projection formula. Only the
fixed kernel in \eqref{eq:kernel} needed an algebraic
representative. The equality therefore holds for arbitrary
topological $F$. Applying the Chern character proves
\eqref{eq:tautological-character}.
\end{proof}

For use as a finite formula, a Mukai vector $(r,c,s)$ has
ordinary character $r+c+(s-r)x_X$. Truncate the exponential
in \eqref{eq:tautological-character} after degree eight,
form \eqref{eq:J}, and multiply its character by
$\exp(-c_1(J_y))$. The Newton identities recover $c_1$ through
$c_4$; higher classes vanish on $Y$. Formula
\eqref{eq:blowup} and the integral bases then give their
coordinates. This covers virtual classes of every rank.
The intermediate arithmetic is rational, but the results are
Chern classes of integral topological indices.

\subsection{Exceptional coordinates}
Applying \eqref{eq:g} to \eqref{eq:curves} gives
\begin{equation}\label{eq:gc}
 gc_0=(5/2,u_0,5/2),\qquad
 gc_2=(-3/2,u_2,-3/2),\qquad q(u_0)=10,\quad q(u_2)=2.
\end{equation}
Substituting \eqref{eq:theta-transport} and \eqref{eq:c2}
in the plane formula gives exactly
\eqref{eq:planes-intro}. For any $a\in\Lambda$,
$\mu(a)^2-\delta\mu(a)$ is even by
\eqref{eq:degree2-products} and the evenness of $q$.
This proves the asserted integrality of both plane coordinates.
Multiplication gives $p_0^2=p_2^2=3$ and $p_0p_2=1$.

Write
\begin{equation}\label{eq:d-ell}
 \begin{gathered}
 \lambda_0=\mu(u_0)+\tfrac52\delta,\qquad
 \lambda_2=\mu(u_2)-\tfrac32\delta,\qquad d_i(a)=q(\lambda_i,a),\\
 \ell_0=\sum_jq(u_0,\alpha_j)R_j-5R_\delta,\qquad
 \ell_2=\sum_jq(u_2,\alpha_j)R_j+3R_\delta.
 \end{gathered}
\end{equation}
The functionals $d_i:R^2\to\Z$ are integral, and $\ell_i$
is their Poincar\'e dual. They are precisely the line
functionals and line classes of the two planes under $\Phi_A$.
The primitive witnesses \eqref{eq:witness-restrictions}
therefore apply to these explicit coordinates.
